\documentclass[10pt,a4paper]{article}

\usepackage[T1]{fontenc}
\usepackage[utf8]{inputenc}
\usepackage{lmodern}
\usepackage{microtype}
\usepackage[a4paper,margin=25mm,headheight=14pt]{geometry}
\usepackage{amsmath,amssymb,amsthm}
\usepackage{booktabs,tabularx,array}
\usepackage{xcolor}
\usepackage{enumitem}
\usepackage{hyperref}

\definecolor{ink}{HTML}{202938}
\definecolor{linkblue}{HTML}{1D4F7A}
\hypersetup{
  colorlinks=true,
  linkcolor=linkblue,
  urlcolor=linkblue,
  citecolor=linkblue,
  pdftitle={Algorithm Frontier: A Curated Catalogue},
  pdfauthor={Algorithm Frontier}
}
\color{ink}
\setlength{\parindent}{0pt}
\setlength{\parskip}{0.45em}
\setlist{nosep,leftmargin=1.3em}
\setlist[description]{style=multiline,leftmargin=28mm,labelindent=0pt}
\renewcommand{\arraystretch}{1.18}
\newcolumntype{Y}{>{\raggedright\arraybackslash}X}
\newcolumntype{P}[1]{>{\raggedright\arraybackslash}p{#1}}

\newcommand{\repo}[2]{\href{https://github.com/icekylinx/algorithm-frontier/#1}{#2}}
\newcommand{\poly}{\operatorname{poly}}
\newcommand{\eps}{\varepsilon}

\newtheoremstyle{recordstyle}
  {0.65em}{0.65em}{\itshape}{}{}{\bfseries}{.5em}{\thmname{#1}\thmnote{ (#3)}}
\theoremstyle{recordstyle}
\newtheorem*{record}{Curated record}

\title{\vspace{-1.2em}\textbf{Algorithm Frontier}\\[0.25em]
  \large A curated catalogue of exact algorithms, tradeoffs, and conditional frontiers}
\author{\normalsize Source-linked research catalogue}
\date{26 October 2026}

\begin{document}
\maketitle
\thispagestyle{empty}

\begin{abstract}
This catalogue collects exact algorithms, approximation schemes, dynamic graph
data structures, and lower-bound records in a common notation.  The entries are
organized by problem family.  Each statement names its input model, output
contract, controlling parameters, and the resource bound that is actually
claimed; incomparable time--space or update--stretch points remain separate.
The repository cards and source maps linked below are the authoritative detailed
records.  This document is a compact mathematical index, rather than a proof
of any result.
\end{abstract}

\section{Scope and notation}
The input-size symbol $H$ denotes the complete binary encoding together with
the polynomial work charged to the stated oracle or primitive.  Polynomial
factors in $H$ are suppressed only when the displayed exponential term is the
point of the record.  The following symbols recur throughout the catalogue:

\begin{center}
\begin{tabularx}{0.96\linewidth}{@{}P{18mm}Y P{18mm}Y@{}}
\toprule
Symbol & Meaning & Symbol & Meaning\\
\midrule
$n,m$ & vertices and edge records & $L$ & binary length of a subset-sum instance\\
$T,Q$ & target budget and support-pattern count & $B,B_0$ & maximum and root quota caps\\
$\Delta,p$ & nonroot incidence and shared-family counts & $U,h$ & logarithmic history and hierarchy scales\\
$r,q,c$ & selected rank, item count, and rejected count & $\eps$ & rational approximation accuracy\\
\bottomrule
\end{tabularx}
\end{center}

A bound on time, writable space, output length, or update work is labelled
explicitly.  A conditional row includes the backend or oracle contract in its
statement.  Lower-bound scope is recorded separately from an upper bound.

\section{Subset Sum}
Both records solve signed binary Subset Sum exactly and return an indexed
witness on YES instances.  Repeated, zero, and negative weights are part of the
input model; $n$ counts weights and $L$ counts the complete encoding.

\begin{record}[faster deterministic point]
For a target $t$, decide whether some subset of the $n$ weights sums to $t$ and
return its original indices.  The reviewed time envelope is
\[
  \poly(L) + \poly(n,L)\,2^{En}, \qquad
  E=\frac{149809813343531}{301754962128739}
    \approx 0.4964618056.
\]
No separate space exponent is claimed.  This point is therefore a time record,
not a simultaneous time--space guarantee.  See the
\repo{problems/subset-sum-fast-deterministic/}{faster-time card} and its
\repo{sources/subset-sum/}{source map}.
\end{record}

\begin{record}[deterministic low-space point]
The same exact decision and indexed-witness task has
\[
  \begin{aligned}
    \text{time}&=\poly(L)\,2^{n/2}, &
    \text{writable space}&=\poly(L)\,2^{E_Sn},\\[-0.15em]
    E_S&=\frac{18261635770550353}{10^{17}}
      =0.18261635770550353.
  \end{aligned}
\]
Read-only input storage is excluded from the writable-space row.  The exact
space exponent supersedes the weaker $137/750$ presentation.  See the
\repo{problems/subset-sum-low-space-deterministic/}{low-space card}.
\end{record}

\section{Numeric targets and quota systems}
These families separate global target budgets from ordered or sparse quota
geometry.  Binary caps are parameters; they are not expanded into unit-sized
copies.

\subsection{Target and support-sensitive polymatroids}
The numeric-target record gives exact target feasibility and scalar optimization,
with approximation and support-sensitive refinements.  In the arbitrary-target
row the envelope is exponential in $T$; fixed row dimension $d$ yields a sharper
stretched-exponential row, and $Q$ distinct eligible-column support patterns
replace an arity restriction in the support-sensitive row.  The common forms
include
\[
  2^{O((T+1)\log(T+2))}\,\poly(H) \quad\text{(arbitrary target)},
  \qquad
  2^{O(T\log(2+Q/T))}\,\poly(H) \quad\text{(support-sensitive row)}.
\]
Each approximation statement carries its own $\eps$ dependence; the source
does not identify approximation with exact decision.  The complete parameter
contract and the two rank-oracle lower-bound scopes are in the
\repo{problems/numeric-target-frontier/}{numeric-target card}.

\subsection{Weighted prefixes, laminar caps, and sparse quotas}
For the ordered quota family, $B$ is the largest target/root cap, $B_0$ is a
distinguished root cap, $\Delta\geq 1$ bounds positive nonroot incidences per
column, and $p$ counts shared weight families.  The reviewed end-to-end rows
include
\[
\begin{array}{rcl@{\qquad}rcl}
\text{one order} &:& 2^{O(B)}\poly(H),
&\text{two orders} &:& 2^{O(B\log(B+2))}\poly(H),\\
\text{$p$ families} &:& 2^{O(pB\log(p+2))}\poly(H),
&\text{laminar root} &:& 2^{O(B^2\log(B+1))}\poly(H).
\end{array}
\]
Sparse mixed quotas retain a separate support-sensitive bound, with the
incidence and root-cap dependence written in the card rather than hidden in
the polynomial factor.  Ground-kernel outputs (packing, coverable-matroid,
transversal, and matching kernels) are separate rows with their own parameter
and output contracts.  See the
\repo{problems/weighted-prefix-quota-frontier/}{weighted-prefix card}.

\section{Dynamic graph algorithms}
The dynamic records preserve the distinction between amortized update work,
stretch, recourse, and query cost.  Here $N$ is a polynomial history cap,
$U=\log N+O(1)$, and $h=\log(U+O(1))$.

\subsection{All-pairs distances and spanners}
For fully dynamic APSP, the equalized point has amortized update work and metric
stretch of the form
\[
  \exp\!\bigl(O(U^{4/5}h^{1/5})\bigr),
  \qquad \text{query time}=\poly(U).
\]
For every fixed positive constant $C$, the second point trades these terms as
\[
  \text{update}=\exp\!\bigl(O(U^{3/4}h^{(C+1)/4})\bigr),
  \qquad
  \text{stretch}=\exp\!\bigl(O(U/h^C)\bigr).
\]
These are incomparable upper-bound points; neither is asserted to be optimal.
The \repo{problems/apsp-frontier/}{APSP card} defines $n,m$ and the positive
polynomial-range edge-length model.

Dynamic spanners retain two constructions because their recourse contracts
differ.  With $m_0$ initial edges and $I,D$ insertions and deletions, the
tracked-terminal row uses an $\exp(O(U^{4/5}h^{1/5}))$ envelope, while the
full-horizon row uses $\exp(O(U^{5/6}h^{1/6}))$.  The latter bounds output
recourse independently of the cumulative insertion count; the former tracks
terminal changes throughout the declared history.  See the
\repo{problems/dynamic-spanners/}{spanner card}.

\subsection{Incremental shortest paths and connectivity}
The incremental directed SSSP record is conditional on a frozen detector and
retirement interface.  It returns a simple path with original edge identifiers,
and its total work is
\[
 O(n+m)+(n+m)\exp\!\left(O\!\left(\log E(N)+
 \sqrt{U\bigl(h+\log(1/\eps)\bigr)}\right)\right),
\]
where the candidate backend has $\log E(N)=O(U^{4/5}h^{1/5})$ and
$0<\eps\leq 1/2$.  Incremental SCC, decremental strong-forest, and dynamic
MSF cards keep their own update model and certificate assumptions; they are
indexed from the \repo{docs/FAMILIES.md}{family guide}.

\section{Matching and graph optimisation}
Exact general-graph matching is a conditional composition.  For $m_s$ retained
edge records after normalization, its backend-dependent time is
\[
 O(n+m)+(m_s+1)\exp\!\left(O\!\left((\log(m_s+2))^{4/5}
 (\log\log(m_s+4))^{1/5}\right)\right).
\]
The output preserves original edge identifiers, including parallel lifetimes.
The \repo{problems/general-matching/}{matching card} states the raw-backend
contract; no new unconditional matching lower bound is claimed.

Degree-constrained subgraphs transfer the same backend to $f$-factors,
interval degree bounds, and cardinality optimization in an explicit loopless
multigraph.  Static and threshold min-cost flow use integral capacities, costs,
and divergence with binary encoding; their raw-cycle lifetime interface yields
\[
 O((N+1)E(S)^4\poly(\log S)),
 \qquad E(S)=\exp\!\left(O((\log S)^{4/5}
 (\log\log(S+4))^{1/5})\right).
\]
Full-flow output is charged separately.  See the
\repo{problems/degree-constrained-subgraphs/}{degree card} and
\repo{problems/min-cost-flow/}{flow card}.

\section{Allocation, robust flow, and scheduling}
The remaining records preserve scenario, selector, and binary-count parameters
because these contracts change the theorem.  Selected splittable flow chooses
$r$ of $q$ commodities, rejects $c=q-r$, and returns a common rate at least
$(1-\eps)$ times optimum.  Its selector base is raised to $c$:
\[
  \Gamma(1/\eps)^c\,\poly(H),
\]
with the selector, routing budgets, and rational-data model stated separately
in the \repo{problems/selected-flow/}{selected-flow card}.  Adaptive robust flow
uses one common selected set and scenario-specific routings; the explicit
scenario list is part of $H$ and exact full-rate admission hardness is kept
separate.

For arbitrary recipient polymatroids, the principal allocation ratio is
\[
 d=4-\frac{3}{r},\qquad \text{value}\geq \frac{\mathrm{OPT}}{d},
 \qquad \text{time}=2^{O(r)}\poly(H),
\]
with an optional $d+\eta$ refinement.  Two-size scheduling has distinct
corank and large-unit branches: factors $3/2$ and $5/3$, respectively, with
times $2^{O(c)}\poly(H)$ and $2^{O(h)}\poly(H)$.  These regimes are not merged.
The \repo{problems/recipient-polymatroid/}{allocation card} and
\repo{problems/two-size-scheduling/}{scheduling card} give the exact oracle,
decomposition, and lower-bound assumptions.

\section{Incoming checkpoint X: draft families}
The complete X increment is indexed here as a research appendix, not as a
promoted catalogue result.  Its 38 numbered payloads are grouped by the
mathematical interface that must be reviewed together.  Support, query time,
and construction time below are candidate quantities; the source cards retain
the exact hypotheses and open proof obligations.

\begin{center}
\begin{tabularx}{0.96\linewidth}{@{}P{38mm}Y P{22mm}@{}}
\toprule
Family & Payload range and candidate frontier & Status\\
\midrule
Exact sensitivity, flow, matching & 01, 09--16; one-failure summaries with
supports such as $O(k^{3/2})$, $O(n^{3/2}\lambda^{1/4})$, and $O(\mu)$ & Draft\\
Subset Sum & 02; conditional deterministic exponent
$158153814247997/318927986919134\approx0.4958919278$ & Draft\\
Spectral, regression, seminorm & 03, 11, 15, 17--19, 26; correlated-failure
feature lifts and weak-coreset transfers & Draft\\
Matroid, cuts, whole blocks & 04--07, 27--29; represented/oracle and exact-label
decoders & Draft\\
Finite groups and metrics & 20, 30--33, 35--38; subgroup, affine, Cayley,
metric, quotient-rank, and obstruction boundaries & Draft\\
Quantum and stabilizer blocks & 21--25, 34; Pauli, qudit, Weyl, and whole-block
models with explicit local dimension & Draft\\
\bottomrule
\end{tabularx}
\end{center}

The X source tree and family-level review cards are source-only.  The proof,
algorithm, and source roles remain open for each family; finite replays check
only bounded identities such as the displayed Subset Sum ledger and the
seven-element non-Fano example.  See
\repo{sources/checkpoint-x/}{the X source map} and the
\repo{reviews/backlog/checkpoint-x-drafts.md}{review backlog}.

\section{Provenance}
Every displayed row points to a problem card containing its hypotheses, proof
map, source map, and review status.  The compact index is
\repo{CONTENTS.md}{CONTENTS.md}; family boundaries and the complete card list
are in \repo{docs/FAMILIES.md}{docs/FAMILIES.md}.  Finite enumeration and
rational-certificate checks are adjunct evidence for selected proof components;
they do not replace a proof of an asymptotic algorithm.

\begin{thebibliography}{9}
\bibitem{catalogue}
Algorithm Frontier repository, \url{https://github.com/icekylinx/algorithm-frontier/}.

\bibitem{families}
Algorithm Frontier, ``Problem families and result cards,''
\repo{docs/FAMILIES.md}{docs/FAMILIES.md}.

\bibitem{complexity}
Algorithm Frontier, ``Complexity notation and parameter policy,''
\repo{docs/COMPLEXITY.md}{docs/COMPLEXITY.md}.

\bibitem{checks}
Algorithm Frontier, verification scripts and finite certificates,
\repo{tools/}{tools/}.
\end{thebibliography}

\vfill
\begin{center}
\footnotesize\textit{The repository cards are authoritative for hypotheses,
proof details, and revisions.}
\end{center}

\end{document}
